%
% Default card theme macros, implemented with tcolorbox.
% Based on https://tex.stackexchange.com/a/185830
%

% Dummy-Contents
\newcommand{\contentA}{}
\newcommand{\contentB}{}

% Old trimmed box config
%\newcommand{\trimmedbox}{\trimbox{0mm 0mm -5mm -5mm}}

% New trimmed box config: more distance between rows of cards
\newcommand{\trimmedbox}{\trimbox{0mm 0mm -5mm -8mm}}

% Card version
\newcommand{\cardversion}{\node at (3,0.5) {v2.0.0};}
%\newcommand{\cardversion}{\node at (3,0.5) {v\input{../.version}};}

% TikZ/PGF settings
%
% Originally using 9cm X 6cm for the total card size.
%
% This sizing is compatible with a regular deck of cards.
%
% This may or may not account for cutting/cropping discreppancies during
% printing.
%
% Alternative: use 8.8cm X 6cm, which may fit better in a card sleeve.
\pgfmathsetmacro{\cardheight}{8.8}
\pgfmathsetmacro{\cardwidth}{6}
%\pgfmathsetmacro{\cardheight}{9}
\pgfmathsetmacro{\imagewidth}{\cardwidth}
\pgfmathsetmacro{\imageheight}{0.75*\cardheight}
\pgfmathsetmacro{\stripwidth}{0.7}
\pgfmathsetmacro{\strippadding}{0.2}
\pgfmathsetmacro{\textpadding}{0.1}
%\pgfmathsetmacro{\titley}{\cardheight-0.1*\strippadding-0.1*\textpadding-0.1*\stripwidth}
\pgfmathsetmacro{\titley}{\cardheight}
\pgfmathsetmacro{\cardpricey}{\cardheight-\strippadding-1.5*\textpadding-0.5*\stripwidth}

% Shapes
\def\shapeCard{(0,0) rectangle (\cardwidth,\cardheight)}
\def\shapeLeftStripLong{(\strippadding,-0.2) rectangle (\strippadding+\stripwidth,\cardheight-\strippadding-\strippadding-1)}
\def\shapeLeftStripBottom{(\strippadding,-0.2) rectangle (\strippadding+\stripwidth,\cardheight/2-\strippadding-\strippadding-1)}
\def\shapeLeftStripShort{(\strippadding,\cardheight-\strippadding-1) rectangle (\strippadding+\stripwidth,\cardheight+0.2)}
\def\shapeRightStripShort{(\cardwidth-\stripwidth-\strippadding,\cardheight-\strippadding-1) rectangle (\cardwidth-\strippadding,\cardheight+0.2)}
\def\shapeTitleArea{(2*\strippadding+\stripwidth,\cardheight-1*\strippadding) rectangle (\cardwidth-2*\strippadding-\stripwidth,\cardheight-2.8*\stripwidth)}
\def\shapeContentArea{(2*\strippadding+\stripwidth,0.5*\cardheight) rectangle (\cardwidth+0.2,-0.2)}

% Stylings
\tikzset{
    % Rounded corners
    cardcorners/.style={
        %rounded corners=0.2cm
    },
    % Element corners
    elementcorners/.style={
        rounded corners=0.1cm
    },
    % Shadow
    stripshadow/.style={
        drop shadow= {
            opacity=.5,
            %shadow,
            color=black
        }
    },
    % Corners and shadow
    strip/.style={
        elementcorners,
        stripshadow
    },
    none/.style={
    },
    % Card image
    cardimage/.style={
        path picture={
            \node[below=-1.5mm] at (0.5*\cardwidth,\cardheight) {
                \includegraphics[width=\imagewidth cm]{#1}
            };
        }
    },
}

% TikZ-Raster
\newcommand{\carddebug}{
    \draw [step=1,help lines] (0,0) grid (\cardwidth,\cardheight);
}

% Card border
\newcommand{\cardborder}{
    \draw[lightgray,cardcorners] \shapeCard;
}

% Card background
\newcommand{\cardbackground}[1]{
    % This draws a border
    %\draw[cardcorners, cardimage=#1] \shapeCard;

    % This does not draw any border
    \path[cardimage=#1] \shapeCard;
}

% Card type
\newcommand{\cardtype}[4]{
    % First we fill the intersecting area
    % The \clip command does not allow options, therefore
    % we have to use a scope to set the even odd rule.
    \begin{scope}[even odd rule]
        % Define a clipping path. All paths outside shapeCard will
        % be cut because the even odd rule is set.
        \clip[cardcorners] \shapeCard;
        % Fill shapeLeftStripLong and shapeLeftStripShort.
        % Since the even odd rule is set, only the card will be filled.
        \fill[strip,#1] \shapeLeftStripLong node[rotate=90,above left=0.9mm,font=\normalsize] {
            \color{white}\uppercase{#2}
        };
        \fill[none,#1] \shapeLeftStripBottom node[rotate=90,above left=0.9mm,font=\normalsize] {
            \color{white}\uppercase{#4}
        };
        \fill[strip,#1] \shapeLeftStripShort;
    \end{scope}

    \node at (\strippadding+\stripwidth-0.28,\cardheight-\strippadding-\strippadding-0.37) {\color{white}#3};
}

% Specific card types
\newcommand{\cardtypeataque}[1]{\cardtype{itembg}{#1}{\hspace{-1mm}\LARGE\bomb}{Ataque}}
\newcommand{\cardtyperecurso}[1]{\cardtype{characterbg}{#1}{\hspace{-1mm}\LARGE\lefthand}{Recurso}}
\newcommand{\cardtypedefesa}[1]{\cardtype{abilitybg}{#1}{\hspace{-1mm}\Large\floweroneright}{Defesa}}
\newcommand{\cardtypeshabilidade}[1]{\cardtype{testbg}{#1}{\hspace{-1.4mm}\huge\ding{78}}{Habilidade}}
\newcommand{\cardtypeattack}[1]{\cardtype{itembg}{#1}{\hspace{-1mm}\LARGE\bomb}{Attack}}
\newcommand{\cardtypeasset}[1]{\cardtype{characterbg}{#1}{\hspace{-1mm}\LARGE\lefthand}{Asset}}
\newcommand{\cardtypedefense}[1]{\cardtype{abilitybg}{#1}{\hspace{-1mm}\Large\floweroneright}{Defense}}
\newcommand{\cardtypeskill}[1]{\cardtype{testbg}{#1}{\hspace{-1.4mm}\huge\ding{78}}{Skill}}

% Apply contour on each word in a way that the whole block gets line breaks
% Thanks to https://tex.stackexchange.com/a/412863
\ExplSyntaxOn
\NewDocumentCommand{\outline}{m}
 {
  \seq_set_split:Nnn \l_tmpa_seq { ~ } { #1 }
  \seq_map_inline:Nn \l_tmpa_seq { \contour{black}{##1} ~ } \unskip
 }
\ExplSyntaxOff

% Card title
\newcommand{\cardtitle}[1]{
    % Title text area
    \node[text width=3.90cm,anchor=north] at (0.5*\cardwidth,\titley) {
        % Within a transparent color box
        \begin{tcolorbox}[beforeafter skip balanced=0,top=0,bottom=0,left=0,right=0,toprule=0,titlerule=0,toptitle=0,bottomtitle=0,bottomrule=0,leftrule=0,rightrule=0,halign=flush center,colback=titlebg,opacityframe=0.6,opacityback=0.6,enhanced jigsaw]
          % Normal
          \color{white}\uppercase{\normalsize #1}

          % With contour
          %\color{white}\uppercase{\normalsize \outline{#1}}
        \end{tcolorbox}
    };
}

% Card content
\newcommand{\cardcontent}[2]{
    \begin{scope}[even odd rule]
        \clip[cardcorners] \shapeCard;
        \fill[elementcorners,contentbg] \shapeContentArea;
    \end{scope}
    \begin{scope}[even odd rule]
        \clip[cardcorners] \shapeCard;
        \node[below right, text width=(\cardwidth-2*\strippadding-\stripwidth-2*\textpadding-0.3)*1cm] at (2*\strippadding+\stripwidth+\textpadding,0.5*\cardheight-\textpadding) {
            %\textit{\glqq\normalsize #1\grqq}
            %\textit{\glqq\small#1\grqq}
            %\textit{normalsize #1}
            \textit{\small#1}
        };
        \node[below right, text width=(\cardwidth-2*\strippadding-\stripwidth-2*\textpadding-0.3)*1cm] at (2*\strippadding+\stripwidth+\textpadding,3) {
            %\vrule width \textwidth height 2pt \\[-2pt]
            \vspace{0.2cm}
            {\scriptsize #2}
        };
    \end{scope}
}

% Card price
\newcommand{\cardprice}[1]{
    \begin{scope}[even odd rule]
        \clip[cardcorners] \shapeCard;
        \fill[strip,pricebg] \shapeRightStripShort;
    \end{scope}
    \node at (\cardwidth-0.5*\stripwidth-\strippadding,\cardpricey-0.1) {\color{}#1};
}

\newcommand{\cutmarksCard}{
}

\newcommand{\cutmarksAThree}{
}

% Cut margins for A4 paper
\newcommand{\cutmarksAFour}{
  % Cut marks for the "0mm 0mm -5mm -5mm" trimmedbox dimensions
  %\begin{tikzpicture}[remember picture, overlay]
  %  % Horizontal marks: left
  %  \draw (-01.20, +19.33) -- (-00.80, +19.33);
  %  \draw (-01.20, +10.33) -- (-00.80, +10.33);
  %  \draw (-01.20, +09.78) -- (-00.80, +09.78);
  %  \draw (-01.20, +00.78) -- (-00.80, +00.78);
  %
  %  % Horizontal marks: right
  %  \draw (+28.60, +19.33) -- (+28.20, +19.33);
  %  \draw (+28.60, +10.33) -- (+28.20, +10.33);
  %  \draw (+28.60, +09.78) -- (+28.20, +09.78);
  %  \draw (+28.60, +00.78) -- (+28.20, +00.78);
  %
  %  %% Vertical marks: bottom
  %  \draw (+00.17, -01.00) -- (+00.17, +00.00);
  %  \draw (+06.18, -01.00) -- (+06.18, +00.00);
  %  \draw (+07.04, -01.00) -- (+07.04, +00.00);
  %  \draw (+13.04, -01.00) -- (+13.04, +00.00);
  %  \draw (+13.90, -01.00) -- (+13.90, +00.00);
  %  \draw (+19.90, -01.00) -- (+19.90, +00.00);
  %  \draw (+20.77, -01.00) -- (+20.77, +00.00);
  %  \draw (+26.77, -01.00) -- (+26.77, +00.00);
  %
  %  %% Vertical marks: top
  %  \draw (+00.17, +21.00) -- (+00.17, +20.00);
  %  \draw (+06.18, +21.00) -- (+06.18, +20.00);
  %  \draw (+07.04, +21.00) -- (+07.04, +20.00);
  %  \draw (+13.04, +21.00) -- (+13.04, +20.00);
  %  \draw (+13.90, +21.00) -- (+13.90, +20.00);
  %  \draw (+19.90, +21.00) -- (+19.90, +20.00);
  %  \draw (+20.77, +21.00) -- (+20.77, +20.00);
  %  \draw (+26.77, +21.00) -- (+26.77, +20.00);
  %\end{tikzpicture}

  % Cut marks for the "0mm 0mm -5mm -8mm" trimmedbox dimensions
  % Difference from the other dimensions: 60x and 30y
  %\begin{tikzpicture}[remember picture, overlay]
  %  % Horizontal marks: left
  %  \draw (-01.20, +19.63) -- (-00.20, +19.63);
  %  \draw (-01.20, +10.63) -- (-00.20, +10.63);
  %  \draw (-01.20, +09.78) -- (-00.20, +09.78);
  %  \draw (-01.20, +00.78) -- (-00.20, +00.78);

  %  % Horizontal marks: right
  %  \draw (+29.20, +19.63) -- (+28.80, +19.63);
  %  \draw (+29.20, +10.63) -- (+28.80, +10.63);
  %  \draw (+29.20, +09.78) -- (+28.80, +09.78);
  %  \draw (+29.20, +00.78) -- (+28.80, +00.78);

  %  %% Vertical marks: bottom
  %  \draw (+00.77, -01.00) -- (+00.77, +00.00);
  %  \draw (+06.78, -01.00) -- (+06.78, +00.00);
  %  \draw (+07.64, -01.00) -- (+07.64, +00.00);
  %  \draw (+13.64, -01.00) -- (+13.64, +00.00);
  %  \draw (+14.50, -01.00) -- (+14.50, +00.00);
  %  \draw (+20.50, -01.00) -- (+20.50, +00.00);
  %  \draw (+21.37, -01.00) -- (+21.37, +00.00);
  %  \draw (+27.37, -01.00) -- (+27.37, +00.00);

  %  %% Vertical marks: top
  %  \draw (+00.77, +21.00) -- (+00.77, +20.00);
  %  \draw (+06.78, +21.00) -- (+06.78, +20.00);
  %  \draw (+07.64, +21.00) -- (+07.64, +20.00);
  %  \draw (+13.64, +21.00) -- (+13.64, +20.00);
  %  \draw (+14.50, +21.00) -- (+14.50, +20.00);
  %  \draw (+20.50, +21.00) -- (+20.50, +20.00);
  %  \draw (+21.37, +21.00) -- (+21.37, +20.00);
  %  \draw (+27.37, +21.00) -- (+27.37, +20.00);
  %\end{tikzpicture}
}
